\scanpage{63}
\mspage{12}
The datum of an $R_i^\omega$, \inserted{which we shall write $R_{\omega i}^{\omega}$, is equivalent to} that of a frame \inserted{$r=(i,\omega)$} of $J$, and defines \uline{two} \deleted{regions} regions $\widetilde D_r$ which contain it, corresponding to the two frames of $(J,\uline{\varepsilon})$ which extend $r$ -- and there are therefore two corresponding regions $\widetilde\Sigma_r$, which are the only ones which contain $R_i^\omega$ as an \uline{interior point} (NB. in the \unsure{standard|standard} region $\widetilde\Sigma_{r_\varepsilon}$, the two other points $R$ apart from $R_0^+$, which is interior, are $R_0^-$ and $R_1^-$, which are not). But if the sphere $\Sigma$ is oriented i.e. if we have given ourselves a marked element \inserted{$\bar\omega_+$} of $\uline{\bar\omega}=\uline{\omega}\wedge\uline{\varepsilon}$, or again an isomorphism $\unsure{\uline{\varepsilon}|\uline{\omega}}\simeq\deleted{\unsure{\uline{\omega}}}\inserted{\uline{\varepsilon}}$, then the point $R_i^\omega$ determines a preferred $\varepsilon$ \inserted{$\in\uline{\varepsilon}$}, by the condition that $\varepsilon$ correspond to $\omega$ under $\uline{\omega}\simeq\uline{\varepsilon}$ or again that $\omega\wedge\varepsilon=\bar\omega_+$, which also means that the orientation $P_+S_0^+R_0^+Q_\infty P_+$ of the region $\deleted{\widetilde\Sigma_r}\widetilde D_r$ is \deleted{that of} $\bar\omega_+$. Thus one can then canonically identify, via $\unsure{\Psi_r|\psi_r}$, $\Pi_1(\widetilde\Sigma,R_i^\omega)$ with $\Pi_{0,3}$, by an isomorphism (which depends on the choice of the orientation $\bar\omega_+$ of $\Sigma$)
\[
(43),\qquad\psi_i^\omega:\Pi_{03}\xrightarrow{\sim}\Pi_1(\widetilde\Sigma,R_i^\omega)
\]
\margin{62}
